% Generated draft: requires mathematical and source review.
\documentclass[11pt]{article}
\usepackage[margin=25mm]{geometry}
\usepackage{fontspec}
\setmainfont{Latin Modern Roman}
\usepackage{amsmath,amssymb,mathtools,unicode-math}
\setmathfont{Latin Modern Math}
\usepackage{enumitem,needspace,xurl,hyperref}
\setlength{\parindent}{0pt}
\setlength{\parskip}{5pt}
\setlength{\emergencystretch}{3em}
\newcommand{\N}{\mathbb{N}}
\newcommand{\Z}{\mathbb{Z}}
\newcommand{\Q}{\mathbb{Q}}
\newcommand{\R}{\mathbb{R}}
\newcommand{\C}{\mathbb{C}}
\newcommand{\F}{\mathbb{F}}
\newcommand{\A}{\mathbb{A}}
\newcommand{\PP}{\mathbb{P}}
\newcommand{\E}{\mathbb{E}}
\newcommand{\eps}{\varepsilon}
\newcommand{\dd}{\,\mathrm{d}}
\begin{document}
\section*{UnsolvedMath ID: 30005983}
\textbf{Source problem code:} catalogue entry
\subsection{Definitions and mathematical statement}
The supplied record poses a construction and comparison problem for \(p\)-energies on spaces without a classical gradient. It does not specify the range of \(p\), a class of underlying spaces, or axioms defining a \(p\)-energy or its canonicity. Consequently, it does not determine a single fully quantified conjecture.

A schematic formulation of the requested comparison is as follows. Once an exponent \(p\), an underlying space with the required additional structure, and a common ambient space \(V\) of functions have been specified, write
\[
\mathcal E_p^{\mathrm{graph}},\qquad
\mathcal E_p^{\mathrm{upper}},\qquad
\mathcal E_p^{\mathrm{KS}}
\]
for candidate energies obtained, respectively, from graph limits, upper gradients, and the Korevaar--Schoen construction. For comparison, if these are nonnegative extended-real-valued functionals, regard each as a functional on \(V\) by assigning \(+\infty\) outside its domain. This is a comparison convention, not an assertion that the source supplies such functionals.

\textbf{Construction and comparison problem.} Determine explicit hypotheses on the exponent, underlying space, auxiliary data, and approximation procedures under which these constructions exist and define a canonical \(p\)-energy. Determine hypotheses under which, with explicitly specified identifications and normalizations,
\[
\mathcal E_p^{\mathrm{graph}}(f)
=
\mathcal E_p^{\mathrm{upper}}(f)
=
\mathcal E_p^{\mathrm{KS}}(f)
\qquad\text{for every }f\in V.
\]
Under the extended-value convention, this equality includes agreement of their finite-energy domains. The source does not specify whether its intended comparison is this exact equality, equality up to a normalization factor, or a weaker equivalence. It also does not define the auxiliary choices from which a canonical energy should be independent. These are unresolved parts of the formulation.

The original supplied excerpt asks how to construct a \(p\)-energy without a classical \(\nabla f\), mentions constructions under different kinds of nonsmoothness, and ends with ``For example:'' without supplying the examples. The explicit demand for canonicity and agreement among the three named constructions occurs in the supplied registry statement; the excerpt alone does not establish that stronger target.
\subsection{Short English statement}
Determine when a space without a classical gradient admits a canonical \(p\)-energy, and when graph-limit, upper-gradient, and Korevaar--Schoen constructions produce the same functional. The exponent range, class of spaces, energy axioms, normalization, and meaning of canonicity remain unspecified. This is a statement-only research-program draft requiring mathematical review.
\subsection{Sources}
\begin{itemize}
\item Supplied local UnsolvedMath registry record, permanent ID \(30005983\), exact source problem code \texttt{catalogue entry}, titled \emph{Construction of p-Energies without Classical Gradients}; collection: Oberwolfach Reports Open Problems. No individual proposer is supplied. The recorded proposed year is 2024.
\item Supplied report attribution: \emph{Recent Developments in Dirichlet Form Theory and Related Fields} (2025), Oberwolfach Reports, DOI \texttt{10.4171/owr/2024/42}, pp.\ 2425--2510. Supplied literature-source label: Oberwolfach source report. \url{https://doi.org/10.4171/owr/2024/42}. The supplied background describes Oberwolfach Reports as documenting mathematical workshops and their open-problem discussions. No precise page or named contribution for this problem is supplied.
\item Kigami, \emph{Conductive homogeneity of compact metric spaces and construction of p-energy}. Supplied literature link: \url{https://ems.press/content/book-files/28253}.
\item Sasaya, \emph{Construction of p-energy measures associated with strongly local p-energy forms}. Supplied literature link: \url{https://doi.org/10.1016/j.jfa.2026.111489}.
\item Literature assessment dated exactly \texttt{2026-08-22}: the supplied curation labels the problem partially solved and states that discrete graph limits, \(p\)-conductive homogeneity, upper gradients, and Korevaar--Schoen constructions provide workable \(p\)-energies under different hypotheses, while no general construction applies to every nonsmooth setting. This is supplied dated curation derived from the OWR record, not an independently verified literature assessment. Its time-sensitive status and cited claims require independent verification.
\item The registry labels the statement \texttt{corrected\_verified}, with the qualification that only formula delimiters, notation, grammar, and typography were repaired and mathematical content and scope were preserved. This supplied editorial label does not establish mathematical correctness or verify the comparison target.
\end{itemize}
\end{document}
